% ==================================================================
%  RelatedWork.tex  --  subsection of Section 1 (Introduction)
% ==================================================================
\subsection{Related Work}
\label{sec:related}

Privacy-preserving clustering has been pursued with both HE and secure
multi-party computation (MPC).

\subsubsection{Averaging-Based Clustering under HE}
K-means and mean-shift share an averaging structure: in every iteration,
a point---a centroid or a shifted sample---is replaced by a weighted
average of the data. Non-interactive clustering under HE has been achieved
for both, for mean-shift~\cite{cheon2019towards} and for
K-means~\cite{yang2024optimized}. The weighted sums inside these averages
require only additions and multiplications, which CKKS evaluates natively;
the remaining non-polynomial steps, division and nearest-center
selection, are approximated by polynomials or reformulated.

Cheon et al.~\cite{cheon2019towards} adapt mean-shift, which seeks the
modes of a kernel density estimator (KDE). In the original algorithm,
every point is shifted toward a mode, and the mode it reaches serves as
its label. Cheon et al.\ instead shift only a small number of points
sampled at random from the data, called dusts, while the KDE is still
evaluated over all points, and they replace the usual Gaussian or
Epanechnikov kernel with a polynomial one. After the dusts have been
shifted toward the modes, a separate point-labeling phase maps every
point to its closest mode. It applies MinIdx, a polynomial approximation
of the position of the minimum, to the squared distances from the point
to all dusts, and multiplies the resulting vector entry by entry by the
number of dusts close to each dust. The label of the point is this
product, a vector over the dusts that is close to $1$ at the dusts
nearest to the point and close to $0$ elsewhere.
The parameters of the scheme, among them the numbers of iterations for
mean-shift and for the polynomial approximations, are chosen separately
for each dataset.

Yang et al.~\cite{yang2024optimized} compute each nearest-centroid
assignment as an encrypted one-hot vector via composite polynomial sign
approximation. They feed this vector directly into a stabilized centroid
update that divides by the public dataset size rather than an encrypted
cluster size. The number of iterations is a constant supplied by the data
owner, with no rule given for choosing it, and the scheme inherits the
limitations of K-means discussed above.

\subsubsection{Privacy-Preserving DBSCAN}
Wang et al.~\cite{wang2022homomorphic} return every comparison to the client. They
encrypt the data under the Brakerski--Gentry--Vaikuntanathan (BGV)
scheme~\cite{brakerski2012leveled}, an HE scheme defined over the
integers modulo a plaintext modulus. To
compare two encrypted values $x$ and $y$, the server homomorphically
subtracts $y$ from $x$, adds a fixed constant~$K$, and sends the resulting
ciphertext $\mathrm{Enc}(x - y + K)$ to the client. The client decrypts it
and returns its most significant bit as the comparison outcome, and the
server drives the clustering on these plaintext outcomes.

In MPC-based schemes, non-colluding servers that hold secret shares
instead evaluate the comparisons jointly. Bozdemir
et al.~\cite{bozdemir2021privacy} execute exact DBSCAN between two
servers, computing distances with arithmetic sharing and performing
comparisons and cluster expansion with Yao's garbled
circuits~\cite{yao1986generate}. The scheme reveals nothing beyond the
output, not even the number or sizes of clusters. A data-oblivious
expansion, whose steps do not depend on the data, requires up to $n$
passes of $O(n^2)$ work; the scheme therefore fixes the number of passes
to a value that sufficed on its test data, or otherwise reveals one bit
per convergence check.

Fu et al.~\cite{fu2024ppa} build on $\rho$-approximate
DBSCAN~\cite{gan2015dbscan}, which clusters the cells of a grid rather
than individual points, and run it between two non-colluding servers.
They construct comparisons from secure most-significant-bit extraction
over additive secret shares. The grid partition reduces the number of
communication rounds to $O(n)$. The scheme, however, is designed for
two-dimensional data, and its output is only guaranteed to lie between
the DBSCAN results for the radii $\varepsilon$ and
$(1+\rho)\varepsilon$.

Both MPC schemes rely on non-collusion among multiple servers and on
server-to-server interaction throughout the computation; as Bozdemir
et al.~\cite{bozdemir2021privacy} themselves note, HE would remove the
non-collusion assumption, at a computational cost.

% [2026-11-11] Summary removed to meet the 13-page limit; original text:
% \subsubsection{Summary}
% Non-interactive clustering under HE thus exists for K-means and
% mean-shift, where each comparison outcome is either consumed by an
% arithmetic update or forms the final label. In DBSCAN, comparison
% outcomes are thresholded again by the $\mathit{MinPts}$ test and then
% chained: two core points linked through a sequence of neighboring core
% points fall into the same cluster, so a single erroneous outcome can merge
% two clusters. Existing DBSCAN schemes have avoided evaluating these
% comparisons in the ciphertext, relocating them to the
% client~\cite{wang2022homomorphic} or to additional
% servers~\cite{bozdemir2021privacy,fu2024ppa,fu2024fss}. We evaluate them
% inside CKKS on a single server. The number of propagation rounds is chosen
% by the client, and we give the client a grid-based estimate that helps it
% choose. The result is a non-interactive DBSCAN in which the server learns
% nothing about the data beyond a small set of public parameters, and whose
% clustering accuracy is on par with plaintext DBSCAN.
